% authority: science-draft
% status: draft/control
% route: candidate-construction-packet
% adoption_status: blocked_adoption_open_continuation
% candidate_constructor_result: constructed_candidate
\documentclass[11pt]{article}
\usepackage[margin=1in]{geometry}
\usepackage[T1]{fontenc}
\usepackage{lmodern}
\usepackage{amsmath,amssymb,amsthm}
\usepackage{booktabs}
\usepackage{longtable}
\usepackage{xurl}
\usepackage[hidelinks]{hyperref}

\newtheorem{definition}{Definition}
\newtheorem{theorem}{Theorem}
\newtheorem{proposition}{Proposition}
\newtheorem{corollary}{Corollary}
\newtheorem{nonconclusion}{Non-conclusion}
\newcommand{\Q}{\mathbb Q}
\newcommand{\One}{\mathbf 1}
\newcommand{\Aspace}{\mathcal A_C}
\newcommand{\Qspace}{\mathcal Q_C}
\newcommand{\Kspace}{\mathcal K_C}
\newcommand{\Sclosure}{S_C}

\title{FiniteSourceClosureConstraintCandidate\_v1\\
P8-T03 Proposal-Only Closure Calculation}
\author{The \AE ther-Flow Research Project}
\date{2026-07-29}

\begin{document}
\sloppy
\maketitle

\section{Decisive result and authority boundary}

This packet constructs one nontrivial finite closure-constraint candidate
from the exact adopted P7 finite source-matter action.  The decisive
Candidate Constructor result is
\[
  \texttt{constructed\_candidate}.
\]
The new object is \texttt{draft/control}, \texttt{proposal-only},
\texttt{source-extension data}, with
\texttt{blocked\_adoption\_open\_continuation}.  It is not a canonical
ontology candidate, an adopted gravitational law, a target-spacetime
effective action, an effective metric, an Einstein--Hilbert action, an
Einstein equation, or an exact-GR recovery theorem.

The construction closes all ten assumptions of
\texttt{LocalEffectiveActionClosureTarget\_v1} only in the exact finite
candidate envelope defined below.  It does not close that target in an
unscoped continuum or physical gravitational regime.  The labels
``gravitational carrier'' and ``closure'' name the proposed P8 role of the
new variable; they do not establish that physical interpretation.

\section{Fixed source basis}

Let \(\Omega\) be one declared finite carrier from the adopted P7-T02
source-matter kernel and let \(P\in\Q_{\geq0}^{\Omega\times\Omega}\) use the
row convention
\[
  \sum_yP_{xy}=1.
\]
The exact P7-T06 construction supplies
\begin{equation}
 C=\frac12(P+P^{\mathsf T}),\qquad
 (L_C)_{xy}=
 \begin{cases}
 -C_{xy},&x\ne y,\\
 \displaystyle\sum_{z\ne x}C_{xz},&x=y,
 \end{cases}
 \label{eq:laplacian}
\end{equation}
and the adopted source-matter action
\begin{equation}
 V_{\mathrm{src}}(u)
 =\frac14\sum_{x,y}C_{xy}(u_y-u_x)^2
 =\frac12u^{\mathsf T}L_Cu.
 \label{eq:source-action}
\end{equation}
P7-T08 gives this exact finite action, its first variation \(L_Cu\), and
its componentwise balance a physical source-matter interpretation by
constitutive postulate in their declared domain.  It does not derive the
P8 object below.

Let \(\Pi_C\) be the connected components of the undirected support
\(\{\{x,y\}:x\ne y,\ C_{xy}>0\}\).  Define
\begin{align}
 \Kspace
   &=\ker L_C
     =\{k\in\Q^\Omega:k\text{ is constant on every }K\in\Pi_C\},\\
 \Qspace
   &=\Q^\Omega/\Kspace,\\
 \Aspace
   &=\Kspace^\perp
     =\left\{a\in\Q^\Omega:
       \sum_{x\in K}a_x=0\text{ for every }K\in\Pi_C\right\}.
 \label{eq:spaces}
\end{align}
The quotient \(\Qspace\) is the proposal-only finite gravitational carrier.
The balanced space \(\Aspace\) is the exact codomain of the adopted
source first variation
\[
  a_C(u):=L_Cu\in\Aspace.
\]
No target atlas, metric, dimension, causal cone, stress-energy tensor, or
benchmark value enters these definitions.

\section{Candidate action family and selected normalization}

\begin{definition}[Finite source closure-constraint family]
For dimensionless coefficients \(\alpha>0\) and \(\beta\in\Q\), define
\begin{equation}
 \Sclosure^{\alpha,\beta}([h];u)
 =\frac{\alpha}{2}h^{\mathsf T}L_Ch
  -\beta\,a_C(u)^{\mathsf T}h,
 \qquad [h]\in\Qspace,\quad u\in\Q^\Omega.
 \label{eq:action}
\end{equation}
Because \(a_C(u)\in\Kspace^\perp\), the right-hand side is independent of
the representative \(h\mapsto h+k\), \(k\in\Kspace\).
\end{definition}

The named candidate
\(\texttt{FiniteSourceClosureConstraintCandidate\_v1}\) is the
source-normalized member
\[
  \alpha=\beta=1.
\]
This normalization is an internal candidate convention: it makes the
solution map equal to the quotient class of the source configuration.  It
is not a derivation or calibration of Newton's constant, the speed of
light, a cosmological coefficient, or any physical unit.  The ratio
\(\beta/\alpha\) remains an explicit free coefficient for any attempted
physical interpretation.

\begin{theorem}[Exact finite variation, solvability, and uniqueness]
\label{thm:closure}
For every \(u\in\Q^\Omega\), the first variation of
Equation~\eqref{eq:action} in the admitted direction
\([\eta]\in\Qspace\) is
\begin{equation}
 \delta\Sclosure^{\alpha,\beta}([h];u)[\eta]
 =\eta^{\mathsf T}\bigl(\alpha L_Ch-\beta a_C(u)\bigr).
 \label{eq:first-variation}
\end{equation}
The Euler--Lagrange constraint is therefore
\begin{equation}
 \alpha L_Ch=\beta a_C(u)=\beta L_Cu.
 \label{eq:constraint}
\end{equation}
It has exactly one solution class in \(\Qspace\),
\begin{equation}
 [h]=\frac{\beta}{\alpha}[u].
 \label{eq:solution}
\end{equation}
For \(\alpha=\beta=1\), the closure map is
\begin{equation}
 \mathcal C_C:\Q^\Omega\longrightarrow\Qspace,
 \qquad
 \mathcal C_C(u)=[L_C^+a_C(u)]=[u],
 \label{eq:closure-map}
\end{equation}
where \(L_C^+\) is the componentwise Moore--Penrose inverse over the
orthogonal complement of \(\Kspace\).
\end{theorem}

\begin{proof}
Symmetry of \(L_C\) gives Equation~\eqref{eq:first-variation}.
Equation~\eqref{eq:constraint} is solvable because its right-hand side lies
in \(\operatorname{im}L_C=\Kspace^\perp\).  Its solution difference obeys
\(L_C(h-\frac{\beta}{\alpha}u)=0\), hence lies in \(\Kspace\).
This proves existence and uniqueness on the quotient and
Equation~\eqref{eq:closure-map}.
\end{proof}

\begin{corollary}[On-shell value]
For the source-normalized member,
\begin{equation}
 \Sclosure([u];u)
 =-\frac12u^{\mathsf T}L_Cu
 =-V_{\mathrm{src}}(u).
 \label{eq:on-shell}
\end{equation}
This algebraic identity is not a physical gravitational-energy statement.
\end{corollary}

\section{Closure-target assumption audit}

The finite candidate instantiates the ten P8-T02 assumptions as follows.
\begin{longtable}{p{0.08\linewidth}p{0.84\linewidth}}
\toprule
ID & Exact finite-candidate instantiation\\
\midrule
A0 &
Inputs are the fixed adopted P7 finite kernel, action, first variation, and
componentwise balance; the new quotient action is explicitly proposal-only.\\
A1 &
\(X=\Omega\); locality is the weighted support of \(C\);
\(\Qspace=\Q^\Omega/\Kspace\); and \(\Aspace=\Kspace^\perp\).  No target
atlas or metric is used.\\
A2 &
The source-to-matter response map is \(u\mapsto a_C(u)=L_Cu\), and the
source-to-closure map is \(u\mapsto[L_C^+a_C(u)]=[u]\).  Domains, codomains,
fibers, and failure branch are explicit.\\
A3 &
Counting measure and weighted radius-one neighborhoods define an exact
finite local operator \(L_C\) of graph order two.\\
A4 &
The regime is the exact declared finite carrier; truncation order is
\(N=2\); the remainder is identically zero.  No continuum scale is inferred.\\
A5 &
The adopted source variation \(a_C(u)\) descends to the
representative-independent response
\(\mathcal E_C([h];u)=\alpha L_Ch-\beta a_C(u)\).\\
A6 &
The linearization \(\alpha L_C\) is symmetric, so the finite Helmholtz
condition holds exactly in the stipulated counting pairing.\\
A7 &
The action is a global finite sum of edge and vertex densities.  A fixed
edge-ownership partition counts every unordered edge once, so overlap
transgressions and the gluing obstruction vanish in this finite analogue.\\
A8 &
Weighted-graph automorphisms act equivariantly.  Component shifts are the
declared gauge relation, and componentwise source balance is exactly the
condition for gauge invariance and solvability.  This is not target covariant
conservation.\\
A9 &
The quotient removes all componentwise zero modes; the finite principal
operator is \(L_C\); admitted variations are quotient classes; and the
closed finite carrier has no additional continuum boundary datum.\\
\bottomrule
\end{longtable}

These entries establish a finite model of the antecedent, not a continuum
local effective gravitational action.  In particular, A4 and A7 are exact
finite specializations, and A8 supplies only a finite balance-to-gauge bridge.

\section{Local density, gluing, symmetry, and boundary controls}

Choose an ordering of the unordered edges
\(E_C=\{\{x,y\}:x<y,\ C_{xy}>0\}\).  The quadratic density
\[
  \ell_{\{x,y\}}(h)=\frac{\alpha}{2}C_{xy}(h_y-h_x)^2
\]
and the vertex density \(-\beta a_x(u)h_x\) sum to
Equation~\eqref{eq:action}.  On a finite cover, assign weights
\(\omega_{i,e}\in\Q_{\geq0}\) with \(\sum_i\omega_{i,e}=1\) for each edge
and similarly for vertices.  Local sums then glue exactly to the displayed
global action; the overlap defect is zero by construction.  This is an
explicit finite bookkeeping theorem, not a theorem about a continuum
variational bicomplex.

Every permutation \(R\) satisfying \(R^{\mathsf T}L_CR=L_C\) acts by
\((h,u)\mapsto(Rh,Ru)\) and preserves \(\Sclosure^{\alpha,\beta}\).
The shift \(h\mapsto h+k\), \(k\in\Kspace\), also preserves the action
because \(a_C(u)\perp\Kspace\).  Conversely, an arbitrary forcing vector
\(a\) yields a representative-independent action and a solvable constraint
if and only if
\[
  \sum_{x\in K}a_x=0
  \quad\text{for each }K\in\Pi_C.
\]
Thus the adopted P7 componentwise balance becomes a typed compatibility
condition; it is not renamed as a Bianchi identity.

\section{Spectral control and fixed two-state calculation}

Let \(\lambda_*(C)\) be the smallest positive eigenvalue of \(L_C\) on
\(\Kspace^\perp\).  For a balanced forcing \(a\), the solution
\([h]=(\beta/\alpha)[L_C^+a]\) obeys
\begin{equation}
 \|[h]\|_2
 \leq
 \frac{|\beta|}{\alpha\lambda_*(C)}\|a\|_2.
 \label{eq:gap-bound}
\end{equation}
This is a finite conditioning bound.  It is not continuum stability,
causal well-posedness, or robustness of a physical metric.

For the fixed P7-T02 forward control,
\[
 P=
 \begin{pmatrix}1/2&1/2\\0&1\end{pmatrix},
 \quad
 L_C=
 \begin{pmatrix}1/4&-1/4\\-1/4&1/4\end{pmatrix},
 \quad
 \lambda_*(C)=1/2.
\]
If \(u=(s,-s)^{\mathsf T}\), then
\[
 a_C(u)=(s/2,-s/2)^{\mathsf T},\qquad
 [h]=[u],\qquad
 \Sclosure([u];u)=-s^2/2.
\]
Equivalently, for an independently supplied balanced
\(a=(s,-s)^{\mathsf T}\), the source-normalized solution is
\(h=(2s,-2s)^{\mathsf T}\) in the zero-mean representative and the on-shell
value is \(-2s^2\).

\section{Coefficient and correction ledger}

\begin{longtable}{p{0.18\linewidth}p{0.18\linewidth}p{0.56\linewidth}}
\toprule
Term & Status & Exact control statement\\
\midrule
\(\alpha\) &
free, then normalized &
Positive quadratic coefficient.  Candidate v1 sets \(\alpha=1\) by an
internal dimensionless normalization, not by target matching.\\
\(\beta\) &
free, then normalized &
Linear source coefficient.  Candidate v1 sets \(\beta=1\) so the closure
map is \([u]\); no physical coupling strength is derived.\\
\(C_{xy}\) &
fixed &
Inherited algebraically from the exact adopted P7 kernel through
\((P+P^{\mathsf T})/2\).\\
Cosmological term &
absent by definition &
No source-derived constant or volume term is supplied.  Absence from this
candidate is not a theorem that it vanishes physically.\\
Higher derivative terms &
absent by definition &
The exact finite action uses the radius-one Laplacian only.\\
Nonlocal terms &
absent by definition &
The candidate includes no pseudoinverse in the action; \(L_C^+\) appears
only in the solved response map.\\
Boundary terms &
zero on declared carrier &
The action is on the full closed finite carrier.  Subset restrictions retain
the exact P7 crossing-edge flux.\\
Remainder &
\(R_{1,2}=0\) &
Exact only for the declared finite candidate.  No continuum expansion or
suppression scale follows.\\
Physical dimensions &
undetermined &
All displayed quantities are dimensionless.  No \(G\), \(c\), length,
energy, or action unit is derived.\\
\bottomrule
\end{longtable}

\section{Failure branches and downstream burden}

The candidate fails closed under each of the following conditions:
\begin{enumerate}
\item if a proposed source \(a\) is not componentwise balanced, the quotient
action is not representative-independent and the constraint is unsolvable;
\item if a proposed coefficient has \(\alpha=0\), there is no invertible
quotient quadratic form; if \(\alpha<0\), the candidate loses its finite
minimum property;
\item if a target atlas, metric, Einstein--Hilbert term, Einstein equation,
or benchmark coefficient is used to interpret or calibrate the construction,
the branch is invalid under the claim boundary;
\item if the carrier is extended beyond the declared finite P7 domain, all
locality, gluing, spectral, and remainder claims must be reproved.
\end{enumerate}

P8-T04 may vary this exact proposal-only candidate and compare the resulting
finite constraint structurally with the field-equation burden.  It must
preserve the decisive limitation: the present equation
\[
  L_Ch=L_Cu
\]
is a finite source closure constraint, not an Einstein-type field equation.
No canonical adoption, effective metric, target stress-energy tensor,
equivalence principle, Einstein equation, exact-GR recovery, benchmark
promotion, proof, publication, push, or completed derivation is authorized.

\begin{nonconclusion}
The candidate constructs a mathematically exact finite variational constraint
and a source-to-closure map.  It does not establish that \([h]\) is physical
gravity, that \(L_C\) is a spacetime differential operator, or that its
componentwise shift identity is diffeomorphism invariance or the Bianchi
identity.
\end{nonconclusion}

\end{document}
